\section{De la Blocklist al Neurosymbolic Hybrid}

El último grupo de slides no se limita a enumerar follow-up papers: responde a la pregunta «¿puede un imperfect judgment model funcionar como componente de un sistema?». Su valor práctico depende de dónde se coloque Delphi, de si se le puede aplicar un override, de cómo exprese la identity y la cultural variation, y de cómo los top-down principles restrinjan las bottom-up predictions.

\subsection{Por qué no basta una Keyword Blocklist}

La slide del Alexa Prize primero celebra al ganador y luego repasa la práctica de seguridad de 2017: usar un conjunto de sensitive keywords para impedir que el sistema hablara de sexo, homosexualidad, raza y otros temas. Una blocklist es fácil de implementar, pero reduce un context complejo a una lista de palabras; deja pasar expresiones dañinas que no contienen las palabras clave y, al mismo tiempo, bloquea por error conversaciones normales, de búsqueda de ayuda o de carácter educativo.

\lecturefigure{slide-45-alexa-prize-sensitive-keywords.jpg}{La sensitive-keyword blocklist del Alexa Prize y sus limitaciones}{Diapositiva V045 recuperada de la grabación oficial; Stanford CS25 Lecture 16}

\teachervoice{La ponente llama a la blocklist el status quo de aquella época y pregunta si en 2021 seguía siendo lo único posible. No sostiene que Delphi pueda sustituir directamente una safety policy, sino que el sistema necesita, como mínimo, entender «por qué las mismas palabras significan algo distinto en situaciones distintas».}

\subsection{Identity-Specific Judgment: más context, también más riesgo}

La identity slide compara el juicio sobre una misma relación o conducta cuando se describe con distintos gender / sexuality. Introducir la identity puede corregir un modelo que trata todo por igual, pero también puede amplificar stereotypes. Al leer la figura conviene comprobar si el model usa de verdad el relevant context o si toma el identity token como shortcut; la evaluación requiere counterfactual pairs y group-level error analysis.

\lecturefigure{slide-46-identity-specific-moral-judgments.jpg}{¿Puede Delphi hacer identity-specific moral judgment?}{Diapositiva V046 recuperada de la grabación oficial; Stanford CS25 Lecture 16}

\subsection{Controllable Generation y Social Norm Alignment}

La parte izquierda de la follow-up slide muestra controllable text generation: dados una persona, una norm o una constraint, se guía al sistema para que genere historias o diálogos más adecuados; el texto de la derecha señala explícitamente que los dialogue systems actuales pueden aceptar a ciegas contenido ofensivo, y que el objetivo es que el agent aprenda la human moral preference. Cuanto más potente es la función, más necesario es informar de quién proviene la preference.

\lecturefigure{slide-47-controllable-generation-bias.jpg}{Delphi como componente commonsense / norm de la controllable generation}{Diapositiva V047 recuperada de la grabación oficial; Stanford CS25 Lecture 16}

La página siguiente muestra la alineación de social norms y values en interactive narratives, y cita investigación sobre agents de videojuegos. El flujo de la figura conecta el narrative state, el value model y la action choice. Respalda el papel de «judgment model como reward / filter» dentro del sistema, pero no garantiza que un long-horizon agent no eluda las restricciones mediante proxy optimization.

\lecturefigure{slide-48-align-social-norms-interactive-narratives.jpg}{Alinear social norms y values en narrativas interactivas}{Diapositiva V048 recuperada de la grabación oficial; Stanford CS25 Lecture 16}

\subsection{Agenda de investigación: Ethically, Socially, Culturally Aware}

La hoja de ruta coloca en el mismo camino tres requisitos, ethically-informed, socially-aware y culturally-inclusive, y fija como next priority la interpretabilidad y la dynamic customization de diverse viewpoints. Las tres palabras no son redundantes: lo ético atiende al harm y a los principles, lo social a las relaciones y las instituciones, y lo cultural a las diferencias de distribución y a quién queda representado.

\lecturefigure{slide-49-ethical-social-cultural-ai.jpg}{La siguiente etapa de la machine ethics: interpretable, dinámica y culturalmente plural}{Diapositiva V049 recuperada de la grabación oficial; Stanford CS25 Lecture 16}

\subsection{Delphi Hybrid: por qué volver a lo Neurosymbolic}

La hybrid title slide anuncia un «Demonstrative Neuro-symbolic Hybrid Moral Reasoning System». Si se vuelve a la estructura de toda la clase, el neural component aporta language y commonsense generation, y el symbolic / principle component aporta top-down constraints. No se trata de poner reglas encima de la salida del modelo, sino de intentar que las intuiciones sobre casos y los principios se verifiquen mutuamente.

\lecturefigure{slide-50-hybrid-moral-reasoning.jpg}{Delphi Hybrid: un neuro-symbolic moral reasoning system demostrativo}{Diapositiva V050 recuperada de la grabación oficial; Stanford CS25 Lecture 16}

La slide de Delphi failures original enumera dos tipos de defecto: lack of commonsense knowledge y lack of interpretability. En un ejemplo, el modelo emite un juicio erróneo ante «genocide if it creates jobs», lo que muestra que la surface association puede imponerse a los principios; además, la falta de explanation dificulta localizar el error. El objetivo del hybrid es añadir a la vez un knowledge path y un constraint path.

\lecturefigure{slide-51-original-delphi-failures.jpg}{Las carencias de sentido común y la falta de interpretabilidad del Delphi original}{Diapositiva V051 recuperada de la grabación oficial; Stanford CS25 Lecture 16}

\subsection{El marco teórico Bottom-Up y Top-Down}

La slide teórica cita el reflective equilibrium de John Rawls: los bottom-up human judgments aportan intuiciones sobre casos, las top-down constraints aportan principios, y ambos extremos deben corregirse una y otra vez. El sistema Delphi original se sitúa sobre todo en el lado bottom-up; la pregunta del hybrid es cómo codificar el moral principle como symbolic constraints que puedan participar en la inference.

\lecturefigure{slide-52-theoretical-framework.jpg}{Reflective equilibrium: bottom-up intuition y top-down constraints}{Diapositiva V052 recuperada de la grabación oficial; Stanford CS25 Lecture 16}

La página siguiente propone como candidatos top-down la common morality de Bernard Gert y sus diez moral rules, por ejemplo no matar, no engañar y cumplir las promesas. La tabla de reglas parece clara, pero sigue necesitando context, exceptions y conflict resolution. El valor didáctico de la figura está en ofrecer un principle vocabulary que pueda auditarse explícitamente, no en afirmar que diez reglas resuelven mecánicamente todos los problemas éticos.

\lecturefigure{slide-53-top-down-common-morality.jpg}{Top-down moral principles: la common morality y sus diez reglas}{Diapositiva V053 recuperada de la grabación oficial; Stanford CS25 Lecture 16}

\subsection{Hybrid Architecture y resolución de restricciones}

La slide de arquitectura detallada encadena la entrada de la situation, la rule retrieval, la commonsense generation, el Maieutic graph y el final judgment. Al leer la figura conviene seguir el flujo de datos de izquierda a derecha: primero se generan candidate reasons, luego se recuperan o activan los principios, después se usa la estructura del grafo para detectar contradictions y, por último, se emiten el judgment y la rationale. El error de cualquier módulo se propaga, por lo que el sistema requiere module-level evaluation.

\lecturefigure{slide-54-hybrid-architecture-detail.jpg}{El pipeline de retrieval, generation, graph y judgment de Delphi Hybrid}{Diapositiva V054 recuperada de la grabación oficial; Stanford CS25 Lecture 16}

La hybrid inference puede representarse con un objetivo simplificado:
\begin{equation}
z^*=\arg\max_z
\left[
\sum_i \alpha_i s_i^{\mathrm{neural}}(z)
+\sum_j \beta_j\mathbf{1}[C_j(z)]
\right],
\end{equation}
donde $s_i^{\mathrm{neural}}$ es la puntuación de evidencia del commonsense / norm model, $C_j$ son las top-down moral constraints, y $\alpha_i$ y $\beta_j$ controlan el peso de ambos tipos de evidencia. La fórmula es solo una abstracción didáctica; un sistema real también debe manejar conflictos entre reglas, fallos de recuperación y la uncertain abstention.

\begin{lstlisting}[caption={Neuro-symbolic moral reasoning 的模块化循环}]
def hybrid_judgment(situation):
    candidate_reasons = commonsense_model.generate(situation)
    principles = retrieve_moral_rules(situation, candidate_reasons)
    graph = build_argument_graph(situation, candidate_reasons, principles)
    judgment, rationale = constrained_inference(graph)
    if low_confidence(judgment) or unresolved_conflict(graph):
        return abstain_with_audit_trace(graph)
    return judgment, rationale
\end{lstlisting}

\subsection{Resultados y limitaciones del Hybrid}

La página «Where does the original Delphi fail?» señala sobre el sistema original las carencias del Maieutic graph y de la architecture, lo que muestra que cada mejora corresponde a una failure conocida y no a la adición arbitraria de módulos. Exige que el evaluator mida por separado la commonsense retrieval, la principle coverage, la graph consistency y el final judgment, para que una puntuación global no oculte el pipeline bottleneck.

\lecturefigure{slide-55-where-delphi-fails.jpg}{Correspondencia entre los módulos del hybrid y los failure modes del Delphi original}{Diapositiva V055 recuperada de la grabación oficial; Stanford CS25 Lecture 16}

Las últimas adversarial result bars comparan el desempeño de Delphi y Delphi-hybrid en queries overall, certain y ambiguous. En lo mostrado en clase, el hybrid obtiene valores más altos, y al lado se indica que combina keywords con interpretable hybrid moral reasoning. Al leer la figura hay que tener en cuenta que el test set sigue siendo limitado: la mejora no demuestra que las top-down rules cubran un mundo abierto, ni permite ignorar la inference latency adicional.

\lecturefigure{slide-56-adversarial-results-hybrid.jpg}{Resultados de Delphi-hybrid en adversarial open-ended moral queries}{Diapositiva V056 recuperada de la grabación oficial; Stanford CS25 Lecture 16}

\begin{warningbox}{Hybrid no significa «neuronal + simbólico = seguridad automática»}
La modularidad mejora la posibilidad de inspección, pero también multiplica las interfaces donde puede haber errores. El retriever puede elegir reglas equivocadas, el generator puede inventar razones y el solver puede llegar a una respuesta consistente a partir de clauses erróneas. El beneficio en seguridad proviene de que el sistema sea observable, descomponible, capaz de abstenerse y revisable por personas, no de la etiqueta de la architecture en sí.
\end{warningbox}

\subsection{Resumen de la sección}

De la keyword blocklist al hybrid, el sistema incorpora progresivamente context, identity, norms y principles. Un componente ético más potente debe aumentar al mismo tiempo la capacidad de auditoría: registrar el origen de los datos, mostrar la rationale provenance, permitir la abstention, admitir el override y dejar la authority dentro de la cadena institucional de responsabilidad humana.
